\begin{definition}
Let $n\ge2$, $q\ge1$ and let $\Omega'\subseteq\R^n$.  The predicate
$\mathrm{CitedCampanatoEmbedding}(n,q,\Omega')$ is the
cited result (C5) of the Note for the formalization (Campanato's theorem,
[G, chapter III]), in the only direction of (4.10) of the thesis that the paper uses:

\emph{if $\Omega'$ is a bounded Lipschitz domain
(\noderef{IsBoundedLipschitzDomain}), then for every $\lambda$ with $n<\lambda<n+q$ and every $N\ge0$ there is
$N_h\ge0$, depending only on $N$, $\lambda$, $q$ and $\Omega'$, such that every
$u$ with $\|u\|_{\mathcal L^{q,\lambda}(\Omega')}\le N$
(\noderef{CampanatoBoundOn}) and $\|u\|_{L^q(\Omega')}\le N$
(\noderef{MorreyBoundOn} with $\lambda=0$) agrees almost everywhere on $\Omega'$
with a function $w\in C^{0,\alpha}(\overline{\Omega'})$,
$\alpha=(\lambda-n)/q$ (\noderef{MemHolderOn}), whose H\"older seminorm on
$\overline{\Omega'}$ is at most $N_h$.}

The quantifier order encodes the dependence stated in the literature: the
H\"older bound is chosen as a function of $\lambda$, $q$, $\Omega'$ and the two
norms $\|u\|_{\mathcal L^{q,\lambda}(\Omega')}$, $\|u\|_{L^q(\Omega')}$ only,
before $u$ itself.

\textbf{On ``Lipschitz domain'', which is an antecedent of the statement and
not a side condition on the parameter.}  (C5) is stated in the paper at l.~264
as ``\emph{if} $\Omega'$ \emph{is a bounded Lipschitz domain}, $n<\lambda<n+p$
and $u\in\mathcal L^{p,\lambda}(\Omega')$, then \ldots''.  That hypothesis is
part of the cited result and is carried here as the first antecedent, through
\noderef{IsBoundedLipschitzDomain}.  Without it the predicate would assert
Campanato's theorem for \emph{arbitrary} $\Omega'$, which is stronger than the
literature and is forbidden by the Note at l.~258 (``stated as in the literature
and not in a stronger form''); the theorem is false on domains with an outward
cusp, where $|\Omega'\cap B_\rho(x)|$ can decay faster than $\rho^n$ and the
Campanato seminorm no longer controls oscillation.  Conversely the predicate is
not vacuous: for a bounded Lipschitz domain the conclusion is
[G, chapter III], and the hypothesis itself is satisfiable --- every open ball of
positive radius satisfies it (\noderef{BallIsLipschitzDomain}), which is also
what the proofs of \S4 and \S5 need in order to apply (C5) at the balls they
construct.

The remaining content of the paper's ``Lipschitz domain with
$\overline{\Omega'}\subseteq\Omega$'' is carried alongside, by the consumers:
$\Omega'$ open, nonempty, bounded, with $\overline{\Omega'}\subseteq\Omega$.
Those clauses predate this predicate and overlap clauses (1)--(3) of
\noderef{IsBoundedLipschitzDomain}; the overlap is a repeated conjunct and is
left in place so that the proof sites discharging them by name continue to
read correctly.
\end{definition}
